\section{\textcolor{Plum}{Methodenkritik} \textcolor{red}{TODO: Titel}}
\label{ch:Diskussion:Methodenkritik}

allgemeine Kritik (bzw. auf beide bezogen)\\
-> Größe der Stichprobe

% - Warum war Methode A effizienter, aber weniger benutzerfreundlich?\\
% - Welche ISO-Kriterien wurden nicht erfüllt – und warum?\\
% - Welche Limitationen haben deine Messmethoden?


% **********************************************************
%           Usability-Test
% **********************************************************
\subsection{\textcolor{Plum}{Usability-Test} \textcolor{red}{TODO: Titel}}
\label{ch:Diskussion:}
...\\

% Kritische Auseinandersetzung mit der verwendeten Methodik zur Modernisierung von Legacy-Systemen

\noindent
- zu wenige Aufgaben -> liegt aber an dem Umfang des Prototypen\\
- Aufgaben nicht relevant genug\\
- Fragen unverständlich formuliert\\
--> Leute mussten zu oft hin und her, könnte die Messdaten verfälscht haben\\

% **********************************************************
%           Nutzerbefragung
% **********************************************************
\subsection{\textcolor{Plum}{Nutzerbefragung} \textcolor{red}{TODO: Titel}}
\label{ch:Diskussion:}
...

% Kritische Auseinandersetzung mit der verwendeten Methodik zur Validierung der gewählten Dashboard-Lösung

% Ein weitere Grund ist, dass ein einseitiger gepaarter t-Test mehr statistische Power haben wird als sein zweiseitiges Gegenstück. In anderen Worten: auch wenn die Wahrscheinlichkeit, einen Fehler erster Art zu begehen, gleich bleibt (das Alphaniveau, das wir gewählt haben), sinkt die Wahrscheinlichkeit, einen Fehler zweiter Art zu begehen.